\chapter{Energy and the Single Invariant}

The car has energy because it is moving --- kinetic energy, what we ordinarily call the energy of motion --- and that energy is, before it is anything else, a square. It is the speed multiplied by itself: a quantity that comes out large whether the car is going forward or back, and never comes out less than nothing. The last chapter left a question hanging --- what single thing all those small leftovers were the shadow of --- and this chapter answers it, and the answer begins here, with a squared energy read before any length is taken, any square root drawn, any ruler laid against the road. The first honest form of the one quantity the chapter is hunting is a square that cannot be negative.

A squared reading has one place it can hide. A record that changes nowhere could still carry one steady value across the whole stretch, and a plain squared energy, blind to a constant, would pass it by. So the boundary is set to read the load: it catches that steady carried value, pins it against a fixed anchor, and makes the squared reading definite instead of merely never-negative. The boundary here is no surrounding world pressing in --- it is a finite account the apparatus keeps at the edge of the stretch, and what it returns is the load the reading must answer for.

Then the chapter does something it could not do with a single reading. It takes two variations --- not one reading set against itself, but two different ways the same reading could have been nudged --- and pairs them, reading the difference they make together. What it is after is the part neither variation can see alone: the cross-term between them, the coupling that shows only when two different nudges are held against each other rather than one against itself.

And here is the landing the whole book has been climbing toward. Clear away the parts a smaller test can already move --- the first-order changes, the ones a steady cruise leaves balanced --- and clear the boundary's load, and exactly one thing is left standing: a single second-order remainder, reached by no first-order test, the one distinction still waiting after every earlier one has been placed. Read it as that coupled remainder, read it as a squared accounting summed over the stretch, read it as the anchored energy-squared --- it is one number under every name. This is the single invariant: the one quantity the book is named for, the name the last chapter's leftover had been waiting for. The chapter does not reach it by importing a length or a limit; it reaches it on a finite stretch, by finite means, and finds that a measured thing comes down, in the end, to one finite quantity.

But that one quantity is not a size alone. Read on the diagonal --- a variation set against itself --- it is a magnitude: how large the thing that survives is, and never less than nothing. Read off the diagonal --- the cross-term, one variation paired against a different one --- the very same number is free to fall below zero. So the single invariant carries two things at once: how large it is, and which way it leans. The leaning is not noise and not a label stamped on afterward; it is read straight off the pairing, as measured as the size beside it, and the apparatus can no more drop it than drop the size. What the leaning comes to mean --- what one direction distinguishes from the other --- is not this chapter's to say; that waits for the chapters ahead. The single invariant the book is named for is a size and a sign together, and everything after this carries both forward, never the size alone.

\section{Energy-Squared Before Norm}

Chapter 7 ended with a precise demand. The finite apparatus had learned how
to hold a support, name its positions, write local equations, choose a current
reading, and carry residual magnitudes through a squeezed tower. Those acts
made the record readable, but they did not yet give the remaining pressure one
quantity. The first answer cannot be a length. It cannot be a completed
geometry. It must be something the finite record can compute before those
later doors open: a nonnegative quadratic read.

Let the supported coordinate readings be
\[
  x_0,\;x_1,\;\ldots,\;x_n
\]
at the named positions of the finite support. The adjacent difference
\(x_{i+1}-x_i\) reads how the carried value changes from one supported name
to the next. This is not a claim that the difference is itself one of the
cubic residuals from the previous chapter; it is a read on the same supported
coordinate record. The energy-squared quantity tanges those adjacent changes,
funging the whole supported record into a single number:
\[
  E^2(x)
    =
  \sum_{i=0}^{n-1} \bigl(x_{i+1}-x_i\bigr)^2 .
\]
Every term in this display belongs to the finite support. The indices run
over adjacent named positions. The differences use only carried coordinate
readings. The summation stops after finitely many terms. Nothing in the
display asks for an outside space, a limiting object, or a square root.

The term \emph{energy} enters here in its most restrained form. It does not
mean force, mass, motion, or a hidden physical substance. It means a carried
quadratic account of change in the supported record. A sharp jump between
adjacent names makes the corresponding square contribute more; a gentle change
contributes less; two adjacent names that carry the same reading contribute
nothing at all. The read therefore turns local adjacent change into a single
finite total without asking for a distance between all possible records.

Nonnegativity comes from the construction itself:
\[
  \bigl(x_{i+1}-x_i\bigr)^2 \ge 0
  \quad\text{for each } i,
  \qquad
  E^2(x) \ge 0 .
\]
The apparatus does not prove this by importing analysis. A square cannot
contribute a negative amount to the finite sum. The energy-squared read is
therefore a size-shaped quantity from the start: it returns a number at least
zero, and it gathers the adjacent changes of the supported record into a
single carried reading.

This is the promised first move after the invariant question. Chapter 7
controlled residual magnitude along a tower. The present section does something
narrower and more direct: it reads one finite coordinate record quadratically. The
reading is founded on the record, not imported into it, because it reuses the
same named support and the same coordinate values that made the preceding
residual talk possible. The apparatus does not ask for trust in a new
background geometry. It asks for inspection of a finite sum of squares.

This also explains why the display uses adjacent differences rather than a new
diagonal term. A term such as \(x_i^2\) would ask the record to measure the
absolute size of each coordinate reading against a privileged zero. This
section is more austere. It measures change across the supported names and
leaves the question of anchoring the whole record for the next finite act.
Energy-squared now reads variation in the supported record; it does not yet
decide which constant records count as rest.

Energy-squared must come before norm because the square root would ask for
more apparatus than this section has earned. A norm would not merely give a
nonnegative number. It would turn that number into a length, compare lengths,
and eventually invite a completed geometry in which limiting readings can be
closed. None of that belongs here. The finite record can already square
adjacent differences and add them. It has not yet earned the right to turn
that quadratic read into a length.

This order matters because a square root can hide a decision. Once the
quadratic read becomes a length, the prose can start to speak as though the
surrounding geometry has already arrived. Measurement refuses that shortcut.
The sum of squares is available before such geometry because it is only a
finite arithmetic event on carried readings. The square root, by contrast,
asks for a richer value discipline and for comparison laws that the present
record has not yet supplied. The section therefore keeps the quadratic
quantity exposed as \(E^2\), not as \(E\).

The distinction also protects the next question. The pure-difference energy
can vanish when all adjacent readings agree:
\[
  E^2(x)=0
  \quad\text{whenever}\quad
  x_{i+1}=x_i \text{ for every } i .
\]
This observation does not settle what makes the reading detect only the
rest record. It deliberately leaves that issue open. If a constant carried
record has no adjacent change, the pure-difference read reports zero. The
apparatus has not failed; it has exposed the exact question the next section
must answer. What finite boundary act, if any, turns zero energy-squared into
genuine rest?

Thus this first section buys one thing and defers several others. It buys a concrete
order-two magnitude the apparatus can compute now:
\[
  x \longmapsto E^2(x).
\]
It defers the square root, the norm, and every completion-facing claim. It
also defers the boundary work that decides when zero energy-squared means
that no activity remains. Most importantly, it does not yet name the final
order-two remainder as the invariant. It only gives the invariant question
its first honest instrument: a finite nonnegative energy-squared reading on
the supported record.

\section[Boundary Accounting]{Boundary Accounting Removes the Vanishing Mode}

The pure-difference energy-squared from the previous section has one honest
blind spot. It reads adjacent change. If the supported record carries the same
reading at every named position, every adjacent difference vanishes. The sum of
squares then reports zero:
\[
  E^2(x)=0
  \quad\text{can still hide a constant carried record.}
\]
This is not a flaw in the construction. It is the construction telling the
truth about its own reach. A read that sees only adjacent change cannot, by
itself, decide whether a constant carried value counts as rest or as an
unaccounted drift. It sees no difference from one supported name to the next.
The record may still carry a value everywhere. Pure difference has exposed the
vanishing mode.

That exposure matters because Measurement does not treat zero as a magical
word. A zero reading means exactly what the instrument has earned. At the
first level, zero energy-squared means that the record shows no adjacent
change across the named support. It does not yet mean that the carried record
has no activity left. The apparatus has to ask a second finite question:
where can a constant carried value go if every interior difference has already
fallen silent?

The answer cannot come from the interior differences alone. Those differences
have already done their work. They compare neighboring names, square the
change, and add the result. If no neighboring pair changes, the interior
account has no further mark to spend. The missing decision lies at the
boundary of the account: the places where the finite support must say what it
accepts as fixed, what it lets pass, and what it refuses to leave floating.

Boundary accounting names that obligation. The boundary does not decorate an
otherwise complete interior read. It belongs to the finite apparatus because
the interior read has an edge. A supported record has first and last names,
allowed entry points, and places where an unresolved carried value can hide
unless the apparatus accounts for them. The boundary tells the record how its
edge participates in the same finite act as the interior differences.

One may call this natural boundary behavior when the boundary terms have no
separate unresolved contribution, but the public point is simpler: boundary
accounting discharges what the interior read cannot decide by adjacency alone.
It says that the edge of the finite support has not been ignored. The
apparatus has checked the terms that would otherwise remain outside the
interior sum. After that check, the constant carried record cannot escape by
claiming that no adjacent pair changed.

This also clarifies why the boundary must do real work. A boundary that merely
names the edge changes nothing; the same constant record would still pass
through the energy-squared read unseen. To earn its place, boundary accounting
must constrain how the record may sit at the edge, tying the carried value to a
finite anchor so that a uniform drift surfaces as a failed accounting rather
than as harmless rest.

The anchor condition performs that finite act. It does not add a second
energy, import a surrounding geometry, or smuggle in a completed norm. It pins
the otherwise free constant carried value at the boundary. Once the anchor is
in place, a record that stays constant across the support cannot keep a
nonzero value everywhere without violating the boundary account. If every
adjacent difference vanishes, the record carries one value throughout the
support. The anchor then funges that single carried value through the fixed
boundary account, and only the zero carried value survives both tests.

The before-and-after contrast is the section's whole mathematical point. Before
the boundary act, the pure-difference read can report zero while a constant
record remains:
\[
  E^2(x)=0
  \quad\text{does not yet force}\quad
  x=0 .
\]
After boundary accounting and the anchor condition, the zero reading gains a
sharper meaning:
\[
  E^2(x)=0
  \quad\Longrightarrow\quad
  x=0 .
\]
The implication is earned. It does not come from the sum of adjacent squares
alone. It comes from the sum of adjacent squares together with the boundary
account that prevents a constant carried record from floating through the
read.

The notation \(x=0\) should be read with the same finite restraint as the
previous section's energy. It does not claim that the apparatus has entered a
completed function space. It says that every carried reading in the finite
record is the zero reading once the boundary account and anchor condition have
done their work. The apparatus has not measured all possible objects. It has
measured this supported record under this boundary discipline.

This is why the anchor does not act as an arbitrary external decree. It
answers the exact ambiguity produced by pure differences. A constant carried
record has no interior strain. If the apparatus wants zero energy-squared to
detect rest, it must decide whether that constant drift belongs to the record
or escapes the measurement. The anchor makes the decision finite: the boundary
allows no hidden constant offset. The record may rest only by agreeing with
that account.

The gain is definiteness at the level now available. Before the anchor, the
energy-squared read was only nonnegative. It could say that the record carried
no adjacent change, but it could not separate zero rest from constant drift.
After the anchor, the same kind of zero reading can identify the zero record.
The apparatus has not changed the arithmetic of the squares. It has changed
what counts as an admissible record for that arithmetic to certify.

This change can be felt in the simplest finite case. Suppose every adjacent
pair agrees. Then the support no longer shows a chain of separate values. It
shows one carried value repeated under several names. The energy-squared read
has no reason to punish repetition, because repetition has no adjacent change.
The boundary account now asks whether that repeated value satisfies the edge
condition. If the anchor pins the edge to zero, the repeated value must be
zero everywhere. The same calculation that once registered only the absence of
adjacent change can now register the absence of any remaining carried value,
because the boundary has removed the place where a uniform drift could hide.

The point is not that the boundary supplies a larger number to add to the old
sum. The point is that the boundary supplies a condition on admissibility. A
record may enter the zero-detection claim only after it has passed the
boundary account. This keeps the arithmetic and the discipline separate. The
arithmetic still says that a sum of adjacent squares vanishes exactly when
each adjacent square vanishes. The discipline says that a record with all
adjacent differences zero must also satisfy the anchor. Together they turn a
flat carried record into the zero carried record.

This separation also prevents the prose from pretending that the anchor is a
physical nail driven into the record. The anchor is a finite requirement on
what the record may count as under measurement. It belongs to the grammar of
admissible readings. The boundary asks for an account; the anchor gives the
account a fixed reference; the energy-squared read then uses the same
pure-difference arithmetic as before. Nothing mystical happens to the sum. The
record simply loses the right to carry an unaccounted constant offset.

That distinction protects the chapter from a common shortcut. One might want
to announce a norm as soon as zero detects only the zero record. Measurement
does not take that step here. Zero detection is necessary for a later length
discipline, but it is not the length discipline itself. The section has earned
an implication about \(E^2\). It has not earned a square root, a metric, a
completed space, or a topology of limiting records.

Nor has the section installed an operator. The finite account has not yet
asked which matrix, stencil, or transformation represents the energy. It has
only said that the boundary and anchor change what the diagonal reading can
certify. The interior differences still provide the energy-squared
contribution. The boundary account decides whether a silent interior hides a
constant carried value. Together they make zero energy-squared detect rest.

The section also does not teach the boundary as an actual load read. That is a
later question. Here the boundary functions as the condition that closes the
vanishing-mode escape. It tells the energy-squared read when zero has become
decisive. The next step will ask how the boundary itself contributes a read in
the account. This section stops before that contribution becomes an object of
measurement in its own right.

The same restraint applies to the coupled-difference material still ahead. The
vanishing mode has been killed, but the chapter has not yet introduced the
coupled cubic difference or the order-two remainder that will eventually carry
the title claim. The present section does not call the boundary anchor the
single invariant. It does not ask the energy-squared read to bear that weight.
It only prepares the finite ground on which the later invariant can be forced
without pretending that a semidefinite read was already enough.

The grammar is therefore exact. First, pure differences make a nonnegative
energy-squared. Second, a constant carried record reveals the vanishing mode.
Third, boundary accounting names the finite obligation at the edge of the
support. Fourth, the anchor condition prevents the constant record from
escaping that obligation. Fifth, zero energy-squared can now detect rest under
that boundary account. Each step answers the previous step's exposed
ambiguity; none imports a completed geometry.

This also makes clear why boundary work belongs
inside the theory. A boundary is not an afterthought added to a finished
calculation. It determines what the calculation can certify. Without the
boundary account, the energy-squared read certifies absence of adjacent
change. With the boundary account and anchor, it certifies absence of the
carried record itself. The same displayed zero changes its authority because
the apparatus has changed what remains admissible at the edge.

That change of authority is the only claim this section needs. It does not say
that every boundary has this effect. It does not say that every anchored
record has already entered a full analytic setting. It says that the finite
Measurement apparatus can turn the pure-difference energy-squared from a
semidefinite read into a zero-detecting read by accounting for the boundary and
pinning the vanishing mode.

With the boundary accounted for and the anchor in place, zero energy-squared
no longer hides a constant carried record; it detects rest. The next section
asks how the boundary itself reads as a load, not merely as the condition that
makes zero detection possible. Yet the gain already recorded here is the deeper
one. A quantity that is only nonnegative can report the size of change; a
quantity that vanishes for the rest record alone can report the absence of
change as a fact about the record itself, not merely about its neighbors. To
earn that second authority with a finite sum of squares and a boundary
account---without a length, a metric, or a completed space---is to learn how
much a measurement can certify before any geometry is invoked to underwrite it.

\section{The Boundary Reads the Load}

The previous section closed one escape. A record with no adjacent change could
still carry one constant value across the whole support. Boundary accounting
and the anchor removed that hiding place, so zero energy-squared could detect
rest under the boundary discipline. That result is not yet the whole boundary
account. It says that the boundary makes the zero test decisive. The present
section asks what the boundary returns when the account reads it.

This question must stay finite. The boundary is not a surrounding analytic
world, and it is not a metaphorical pressure pressing on the record from
outside. It is a named place in the finite account where a carried record has a
readable value. When the record is carried all the way to that place, the
account receives a loaded boundary value. To call the value \emph{loaded} means
that it is not decorative: it is the value at the boundary position that the
account must compare with the anchor.

The distinction is small but important. A boundary condition says what must be
true for the record to count as admissible. A boundary reading says what finite
value the account obtains at the boundary. The former closes the door on a
hidden constant drift. The latter shows the reader where the closing decision
is read. Measurement needs both. Without the condition, the value has no
discipline. Without the reading, the condition can feel like a rule announced
after the calculation rather than a value the account can inspect.

The reading itself is a finite selection. With the supported values
\(x_0,\ldots,x_n\) funged to the load boundary, the loaded boundary value is
the value standing at that named position,
\[
  v_{\mathrm{load}}(x) = x_n ,
\]
a single carried reading chosen at a finite name---no sum, no length, and no
appeal to a record outside the support.

For a constant carried record, the loaded boundary value has a particularly
plain role. If the record has the same value at every supported name, then the
value read at the load boundary is not a new kind of value. It is the same
carried value reached at the boundary. The interior has already gone silent;
there is no adjacent change left to spend. The boundary read therefore becomes
the finite place where the remaining constant value must answer.

This is why the section can speak of a boundary load without turning load into
an image. The boundary load is the carried value as it is read at the boundary
under the load account. It is not weight, force, cost, or burden. It is a
boundary reading. The account has a finite record, a boundary at which that
record can be read, and an anchor against which the reading is compared.

The constant case gives the clear chain:
\[
  x = v_{\mathrm{load}}(x),\qquad
  v_{\mathrm{load}}(x) = b_{\mathrm{load}},\qquad
  b_{\mathrm{load}} = 0,
  \quad\text{so}\quad x = 0 .
\]
Here \(x\) is the constant carried record. The term
\(v_{\mathrm{load}}(x)\) names the value of that record at the load boundary.
The term \(b_{\mathrm{load}}\) names the anchored boundary load. The final zero
is the anchored value under the boundary account. The display does not ask for
a surrounding space or a completed length. It only follows the finite read
from the carried record to the boundary and from the boundary to the anchor.

Each equality in the chain has its own work. The first equality says that a
constant carried record is already determined by its loaded boundary value. If
the same value is carried throughout the support, then reaching the load
boundary does not change what is being carried. The boundary reads the value
that the flat record has everywhere.

The second equality says that the loaded boundary value is the boundary load
under the anchored account. This is the point at which the boundary reading
stops being a free end. The account does not merely discover some value at the
edge and leave it unjudged. It compares that value with the anchored boundary
load. The anchor gives the boundary reading a fixed reference.

The third equality says that the anchored boundary load is zero. This is not a
new energy calculation. It is the finite boundary answer. Once the record has
been carried to the boundary and the boundary value has been identified with
the anchored load, the anchor returns the zero value. The constant carried
record therefore cannot remain nonzero while also satisfying the load read.

The conclusion \(x=0\) is then not an extra decree. It is the end of the
chain. A constant record equals its loaded boundary value; the loaded boundary
value equals the anchored boundary load; the anchored boundary load is zero.
Therefore the constant record is zero. This is the same zero-detection gain
from the previous section, but now the boundary contribution has become an
explicit read rather than only a condition on admissibility.

The order of that chain should not be compressed too quickly. If the account
begins with \(b_{\mathrm{load}}=0\), the reader sees only the anchor and may
miss how the carried record reaches it. If the account begins and ends with
\(x=0\), the boundary has disappeared again into the conclusion. The useful
middle term is \(v_{\mathrm{load}}(x)\). It marks the finite read that carries
the record to the boundary without changing the record into a new object.

That middle term also explains why the boundary read belongs after boundary
accounting rather than before it. Before the vanishing mode has been exposed,
there is no reason to ask the boundary to answer for a constant carried record.
The pure-difference read first shows that a flat record can pass through the
interior account unseen. Boundary accounting then says that such a record must
answer at the edge. Only after those two moves does the loaded boundary value
become the right object of attention.

The loaded boundary value is therefore not an optional ornament on the anchor.
It is the finite passage between the flat record and the anchored answer. The
record reaches the boundary with one carried value. The boundary reading names
that value in the load position. The anchor then supplies the fixed zero answer
for that position. The account has not changed the record by hand; it has
funged the remaining constant value through the boundary read.

The section's language should be heard with that restraint. To say that the
boundary \emph{reads} the load means that the finite boundary account returns a
value. It does not mean that the boundary chooses, acts, or judges. The account
does the reading because the apparatus has named the support, the boundary, and
the anchor. The boundary load is the value the account tanges at that place
under that discipline.

This also clarifies why the boundary read is not simply another interior
difference. Interior differences compare neighboring supported names. They
are excellent at seeing change along the support. They are silent on a record
that carries one value everywhere. The loaded boundary value answers a
different question: when the record reaches the boundary, what value does it
present to the anchor? That question cannot be answered by adding another
interior adjacent difference, because the problem is precisely the absence of
interior change.

Nor is the loaded boundary value a second energy. The energy-squared read
still does the adjacent-difference accounting. The boundary read supplies the
finite value that lets the remaining constant mode answer to the anchor. The
two acts are complementary at this level. One reads change across the support.
The other reads the carried value at the boundary. Together they explain why a
zero energy-squared record under the boundary account no longer conceals a
uniform drift.

The load read also protects the reader from a false shortcut. It would be easy
to say that the anchor kills the constant record and move on. That sentence is
correct as a compressed result, but it hides the finite route. The constant
record must first be seen as the value it presents at the load boundary. Only
then can the anchor turn that boundary value into zero. The route matters
because Measurement does not let a conclusion float free of the read that
earns it.

The route also protects the public vocabulary. The phrase \emph{loaded boundary
value} should carry most of the work, because it says all three things at once:
the value is loaded, it is read at the boundary, and it remains a value in the
finite account. The shorter \emph{boundary load} can stand in once the
discipline is clear, but it should never float as a bare burden. The load is
not a metaphor for weight. It is the boundary value under the anchor-facing
read.

This matters because the same ordinary word appeared earlier in the book under
different pressure. Here it has a narrower job. The load is not the burden of
standing in a support, and it is not a cost paid by the apparatus. It is the
specific boundary reading that receives the constant carried value and submits
it to the anchor. The phrase must keep pointing back to the boundary read, or
the section will blur a measured value into a general metaphor.

The same restraint keeps the section from overextending. The boundary read does
not create a length discipline. It does not complete the record into a larger
setting. It does not install a device for representing the energy. It only
gives the boundary account a value that can be compared with the anchor. The
chapter still has more work to do before the later order-two claim can be
forced.

What the section buys is therefore exact and limited. It turns boundary
accounting from a closing discipline into an explicit finite read. It tells the
reader how a constant carried record reaches the boundary, how the loaded
boundary value is identified with the anchored boundary load, and why that
anchored load returns zero. The zero record is not reached by assertion. It is
reached by following the finite boundary reading.

What the section does not buy is just as important. It does not compare two
later readings. It does not introduce the next comparison. It does not claim
that the chapter's title has already been won. It does not replace the
energy-squared read with a length or an installed representation. The boundary
has become readable as a load, but the later comparison work still has not
begun.

The final order of the argument is now visible. Energy-squared first reads
adjacent change. Boundary accounting then prevents zero energy-squared from
hiding a constant carried record. The present section opens the boundary
account and shows the finite load reading inside it. A constant carried record
equals the value it presents at the load boundary, that boundary value reaches
the anchored load, and the anchored load is zero.

Nothing further is needed for this section's claim. The reader can now see why
the previous implication was earned rather than imposed. If the energy-squared
read reports zero, the constant case is the remaining danger. If the constant
case reaches the load boundary, it presents the loaded boundary value. If that
value reaches the anchored boundary load, the anchor returns zero. The
remaining constant value has nowhere else in this finite account to hide.

That is the exact stopping point. The boundary account has become an explicit
finite loaded boundary reading, and for a constant carried record that reading
reaches the anchored zero value. The next work cannot be another repetition of
the zero test; it must begin only when the chapter is ready to compare the next
finite readings without pretending that the comparison has already been made.
What the section secures, though, is a quiet principle the rest of the chapter
will rely on: a boundary earns its authority not by decree but by returning a
value the account can read, so that the vanishing of a record is something
inspected at a named place rather than asserted from outside. A conclusion that
can be read is a conclusion that cannot hide its grounds.

\section{The Coupled Difference}

The next comparison begins only after the preceding debts have been paid. One
side may read its residual. The other side may read its residual. If either
residual still speaks at first order, then the comparison has not yet reached
the invariant question. It has only found a remaining local imbalance. The
chapter therefore asks for a coupled difference: a finite read that tanges the
two residual lines while separating them from the part that neither side
can see alone.

Write the two first-order residual reads as
\(r_{\mathrm{L}}\) and \(r_{\mathrm{R}}\). For a left variation
\(v_{\mathrm{L}}\) and a right variation \(v_{\mathrm{R}}\), the coupled
difference has the form
\[
  \Delta_{\mathrm{coupled}}(v_{\mathrm{L}},v_{\mathrm{R}})
    =
  r_{\mathrm{L}}(v_{\mathrm{L}})
  + r_{\mathrm{R}}(v_{\mathrm{R}})
  + \delta^2(v_{\mathrm{L}},v_{\mathrm{R}}).
\]
The first two terms record what each side can still report on its own. The
last term records the order-two remainder that appears only when the two reads
are funged together. It is not a continuum limit, and it is not a completed
function-space norm. It is the finite remainder left after the boundary account
has already removed the false zero mode and the two first-order residuals have
been exposed.

This is why the residuals must vanish before the chapter names the invariant.
If the left residual remains, the coupled difference still reports a left
failure. If the right residual remains, it still reports a right failure. Only
when both first-order reads have cleared does the coupled difference become the
second variation itself:
\[
  r_{\mathrm{L}}(v_{\mathrm{L}})
  =
  r_{\mathrm{R}}(v_{\mathrm{R}})
  =
  0
  \quad\Longrightarrow\quad
  \Delta_{\mathrm{coupled}}(v_{\mathrm{L}},v_{\mathrm{R}})
    =
  \delta^2(v_{\mathrm{L}},v_{\mathrm{R}}).
\]
The implication does not discard the residual reads. It funges them forward as
the finite witnesses that the first-order debts have ended. Once they end, the
comparison has exactly one remaining value to read.

That remaining value is the coupled second variation. It is \emph{second} not
because the chapter has smuggled in a calculus of completed spaces, but because
the one-sided residuals have already taken the first-order positions. It is
\emph{coupled} because no single side owns it: the value belongs to the paired
read that survives after each side has said everything it can say alone. This is
the precise sense in which the invariant is single---not the larger residual,
not the sum of the two, but the one finite remainder that lives only in the
pairing, the value still to be measured once every first-order account has
fallen silent.

\section{Second Variation as the Single Invariant}

On the diagonal, the paired read becomes the anchored quadratic read. Write
\(\delta^2(v)\) for the coupled remainder when the same finite variation is
read against itself. Write \(B(v,v)\) for the same order-two value when the
chapter emphasizes its quadratic accounting. Write
\(E^2_{\mathrm{anchored}}(v)\) when the chapter emphasizes the energy-squared
read after the boundary account has supplied the anchor-facing load. These are
not three competing quantities. They are three names for one finite read:
\[
  I(v)
    =
  \delta^2(v)
    =
  B(v,v)
    =
  E^2_{\mathrm{anchored}}(v).
\]

Read as a sorting, the identity is the degenerate pigeonhole. A measurement
places one distinction into one record-slot, and the chapter has worked by
clearing, order by order, every distinction a lower test can place: the
first-order residuals are sorted into their holes and vanish. One distinction
outlasts that sorting, reached by no first-order test---a single order-two
remainder, the one pigeon that remains. The three names on the left count that
one pigeon, not three; \(I(v)\) is \emph{single} because, once every
first-order placement has been made, exactly one distinction is still waiting
and exactly one slot is left to hold it.

This identity is the promised landing point of the chapter. Energy-squared
comes before any length claim. Boundary accounting prevents the zero
energy-squared read from hiding a constant carried record. The loaded boundary
read submits that remaining constant value to the anchor. After those
first-order and boundary debts clear, the coupled difference leaves a single
order-two remainder. The chapter names that remainder \(I(v)\).

The name \emph{single invariant} therefore carries a narrow burden. It does not
assert that every later apparatus is already in place. It says that, at
this finite stage, every permitted way of reading the surviving order-two value
returns the same anchored quadratic quantity. The mixed read, the second
variation, the diagonal quadratic read, and the anchored energy-squared read
all point to one value once the residual terms have vanished.

The coincidence is not a notational accident; it is forced. The term
\(\delta^2(v)\) is the coupled remainder read against a single variation rather
than a pair, so the off-diagonal coupling collapses onto one finite quadratic.
The term \(B(v,v)\) is that same quadratic seen as an accounting: the order-two
contributions gathered and summed over the finite support, the very bookkeeping
that opened the chapter. And \(E^2_{\mathrm{anchored}}(v)\) is that bookkeeping
once the boundary account has pinned the constant mode, so the read is definite
rather than merely nonnegative. Three vantage points---coupling, accounting,
anchored energy---funge into one finite read. The apparatus tanges that
surviving order-two value a single time; the three names are only the angles
from which it is seen.

This also fixes the order of later questions. The invariant has a name, but
the chapter has not yet supplied the operator that carries it through every
comparison. It has not yet supplied quantitative control for all later finite
families. It has not taken a square root, introduced a norm, or opened a
completed-space discipline. Those remain honest debts. What has landed here is
more modest and more precise: the finite identity that lets the next chapter
begin from a named invariant rather than from another zero test.

That modesty is the measure of what has been earned. Reduce a measurement to
the one pressure that outlasts every first-order account and every boundary
debt, and what remains is a single nonnegative number---the same number whether
it is read as a coupled remainder, as a quadratic accounting, or as an anchored
energy-squared. No length was taken, no norm was assumed, and no continuum was
ever invoked; the value stands on a finite support and is reached by finite
means. That a measured thing should come down, in the end, to one finite
quantity, read the same way under every name it answers to, is the claim the
book's title makes and the place the whole book has been climbing toward.

One reading of that single value the chapter has held back, and it is the reading the rest of the book most needs. The number just named is the invariant read on the diagonal --- a variation set against itself --- and read that way it is a size: nonnegative, the magnitude of the surviving order-two pressure. But the diagonal is only one way to read the coupled remainder. Read off the diagonal, as the cross-term that pairs one variation against a different one, the same invariant is no longer bound to stay nonnegative; the pairing that gave the remainder its meaning can return a value below zero as readily as above it. The diagonal reads the invariant's size; the off-diagonal reads what the size alone cannot carry --- a sign.

So the single invariant is not a magnitude only. It carries two finite aspects at once: how large the surviving pressure is, and which way it leans. The sign is not noise left over from the construction, and it is not a convention the apparatus stamps on afterward; it is read straight off the cross-term, the same finite remainder the chapter has been earning, now seen in the one respect the diagonal hid. A measurement reduced to one quantity is reduced to one quantity with a direction, and the apparatus can no more drop the direction than the size without ceasing to report what it measured.

What that sign comes to mean is not this chapter's to say. Here the apparatus has established only that the single invariant has a sign to carry --- a finite orientation read off the pairing, as measured as the magnitude beside it. Later chapters will take up what the orientation distinguishes, and turn the bare fact that the remainder can lean one way or the other into a reading the apparatus can name. The chapter closes one step more honestly for it: the single invariant the book is named for is a size and a sign together, and the geometry that follows funges both forward, never the magnitude alone.

Everything before this identity was a way of earning the right to name the
single invariant; everything after it is a way of funging that one value
forward.

\section*{Coda: Energy in Review}

Chapter 8 took the last chapter's leftover --- the one thing all those small leftovers were the shadow of --- and found it. It got there by reading an energy as a square before any length; by setting the boundary to catch and pin a steady carried value; and by pairing two variations to read the part neither sees alone. When every first-order account had been cleared, one thing was left standing, read the same way under every name it answered to. That is the chapter in review. But set down inside that finding, in the coda's own plain register, is an eighth part of the thing the chapters have been assembling.

The seven parts already on the bench were each something the value could do or be --- differ from place to place, move, be read through a floor of disturbance, take a size on a finite frame, be decided, hold steady, live on a chosen field. The eighth is different in kind. It is not another thing the value can do; it is the one quantity that outlasts everything the value does --- strip away every change a smaller test can reach, and this is what remains, the same number however it is read. And it does not arrive as a size alone: it carries how large it is and which way it leans, the leaning read straight off pairing the value's variations against one another. That is the eighth part. Bare, as the others were: not what the leaning means, not what it is for --- only that the value has one quantity that outlasts its every variation, carrying a size and a way it leans.

What the leaning means does not matter right now. Like the seven before it, the eighth is set down and left unexplained; what it is for waits on the parts not yet handed over. The book keeps its one habit --- one piece per chapter, in its turn, none before the construction reaches it. Eight pieces now lie in order: a value spread across places, the same value moving, the value read through a floor of disturbance, the value sized on a finite frame, the value decided, the value at rest, the value living on a chosen field, and now the one quantity that outlasts every variation of it. A reader who has watched all eight go down may feel the shape pulling tighter, though the book still will not say it.

For now there are eight parts, and the eighth is as bare as the rest. It names no leaning in particular, fixes no size, and leans on nothing still to come: a value that can differ from place to place, a value that can differ from before to after, a value read through a floor of disturbance, a value with a size on a finite frame, a value committed to a definite reading, a value that can hold steady, a value that lives on a chosen field of places, and now the one quantity that outlasts its every variation, carrying a size and a way it leans. It is an eighth mark of a second kind, set down beside the others and waiting, as they wait, for what has not been handed over to give it its place.
